\section{Words and strict drift}\label{word:section}

Fix $q>4$, and put
\begin{equation}\label{word:parameters}
 t=\sqrt q,\qquad p=\frac{t}{t+2},\qquad u=\frac{1-p}{16}.
\end{equation}
Thus $1/2<p<1$. For a rooted connected planar map $M$, define
\begin{equation}\label{word:weight}
 a(M)=\sum_{A\subseteq E(M)}
       t^{\,2k_M(A)+|A|-|V(M)|-1}.
\end{equation}
The decoration-summed weight in \eqref{eq:finite-law} is $t\,a(M)$. Since this
factor is independent of $M$, the marginal law of $M_n$ assigns
probability proportional to $a(M)$ among rooted maps with $n$ edges.
We also adjoin a trivial one-vertex, zero-edge map, whose weight is $1$,
and set
\[
 a_n=\sum_{\substack{M\text{ rooted planar map}\\ |E(M)|=n}}a(M),
 \qquad a_0=1.
\]
A directed root is a distinguished dart; in particular, a loop has two
darts. Rooted maps are counted up to root-preserving,
orientation-preserving isomorphism.

\subsection{Reduction of words}

Consider independent letters from the alphabet
\[
 \mathtt b_1,\ \mathtt b_2,\ \mathtt o_1,\ \mathtt o_2,\ \mathtt F
\]
with probabilities
\begin{equation}\label{word:letterlaw}
 \P(\mathtt b_1)=\P(\mathtt b_2)=\frac14,
 \qquad
 \P(\mathtt o_1)=\P(\mathtt o_2)=\frac{1-p}{4},
 \qquad \P(\mathtt F)=\frac p2.
\end{equation}
The letters $\mathtt b_i$ create burgers of type $i$. Read a word from
left to right. An order $\mathtt o_i$ removes the freshest remaining
burger of type $i$, and a flexible order $\mathtt F$ removes the freshest
remaining burger of either type. An order with no eligible burger is
unfilled; it is never served by a later burger.

The \emph{reduction} of a finite word is its list of unfilled orders,
followed by its list of surviving burgers, with each list retaining its
original chronology. A word is \emph{empty-reducing} if both lists are
empty. Throughout, a list of burgers is ordered from oldest to freshest.

\begin{lemma}[Reduction above an older supply]\label{word:reduction}
Feed a finite word an arbitrary older list of burgers. Every match
between two letters of the word is the same as when the word is read
from an empty supply. The orders unfilled in the latter evaluation act
on the older list in their original chronology, and the surviving burgers
of the word remain above the surviving older burgers. In particular,
finite words may be replaced by their reductions when concatenating
them. Once a suffix has a surviving burger, its freshest survivor is
unchanged by prepending further letters.
\end{lemma}

\begin{proof}
Call burgers created within the word new. At each order, an eligible new
burger is fresher than every old burger and therefore has priority.
Induction over the letters shows that the available new burgers evolve
exactly as they do with empty initial supply. Only orders for which no
eligible new burger exists can affect the old supply, and they do so in
their original order. The new survivors lie above all remaining old
burgers, which proves the first assertion and the concatenation rule.
To obtain the last assertion, first evaluate the prepended word and then
apply the rule to the original suffix. Its surviving burgers remain
unchanged and lie above any extra surviving supply.
\end{proof}

\subsection{An exact counting identity}

The tree-rooted-map encoding goes back to Mullin~\cite{Mullin1967};
Bernardi gives its two-parenthesis formulation~\cite{Bernardi2007}.
The inventory encoding with flexible orders is due to
Sheffield~\cite[Section~4]{Sheffield2016}.
We give the finite weighted identity in the form needed here. Its proof
uses a spanning-tree expansion determined by the embedding; related
embedding-based activity expansions of the Tutte polynomial were
introduced by Bernardi~\cite{Bernardi2008}. The activity convention below
is fixed explicitly, and the identity is proved directly.

\begin{proposition}[Finite word identity]\label{word:identity}
For every integer $n\ge0$, an independent word of length $2n$ satisfies
\begin{equation}\label{word:identity-equation}
 \P(\text{the word is empty-reducing})=a_nu^n.
\end{equation}
\end{proposition}

\begin{proof}
We give the bijection and the weight calculation, including their rooting
conventions.

\paragraph{Typed words and spanning trees.}
In an empty-reducing word, replace each flexible order by the fixed type
of the burger it consumes. The resulting \emph{identified word} has two
types of opening symbols, the burgers, and matching closing symbols,
the orders. Each type separately is a Dyck word: its running balance is
nonnegative, its final balance is zero, and matching follows its stack.
Thus an identified word is a shuffle of two Dyck words.

Such shuffles encode rooted planar maps with a distinguished spanning
tree $\mathcal T$. To see this, walk around a thickening of $\mathcal T$
in the orientation of the sphere, starting with the root dart. At a tree
dart, cross its edge and continue with the successor of the opposite dart
in its vertex's cyclic order. At a non-tree dart, continue with its
successor at the same vertex. This tour visits every dart exactly once.
The first and second occurrences of a tree edge give a type-1 opening
and closing; those of a non-tree edge give a type-2 opening and closing.
The tree pairs are noncrossing in contour order. The non-tree pairs are
also noncrossing, because their edges are disjoint arcs in the disk
complementary to the thickened tree.

Conversely, the type-1 symbols construct a plane tree by descending a new
edge at an opening and returning at its closing. Place the type-2 stubs
between successive tree steps in the given order, and join their matched
pairs in the complementary disk. Noncrossing matching determines these
arcs up to an orientation-preserving homeomorphism, so this constructs a
unique planar map. The first symbol specifies its root dart. The same
construction applies when there are no tree edges: all the stubs lie at
one vertex. At the end/start seam, retain the indicated cyclic order of
stubs. More explicitly, if $\alpha$ pairs matching positions and $\tau$
is the cyclic successor of positions, the reconstructed vertex successor
is $\tau\circ\alpha$ on tree darts and $\tau$ on non-tree darts.
This verifies that the two constructions are inverse, including the
rooting convention.

\paragraph{Flexible marks and activity.}
Say that two edges cross in the tour order when their occurrences
alternate. Such edges necessarily have opposite types. Call an edge
\emph{active} when it crosses no edge with a later first occurrence.
For an edge whose occurrences are $a<b$, its closing order can be made
flexible precisely when its burger is freshest among all burgers then
present. The obstruction is an edge with occurrences $c<d$ satisfying
$a<c<b<d$. Thus the eligible closings are exactly the active edges.
Any subset of them can be marked flexible while retaining all matches:
an induction through the word verifies that each changed order still
consumes the same burger. Conversely, every flexible mark in a word must
be eligible in its identified word.

Every fixed identified word with $n$ matched pairs has probability
$u^n$, since a burger and a fixed order contribute
$\frac14\frac{1-p}{4}=u$. Replacing a fixed order by a flexible one
multiplies its probability by
\[
 \frac{p/2}{(1-p)/4}=\frac{2p}{1-p}=t.
\]
If $b(M,\mathcal T)$ denotes the number of active edges, the total
probability of all its possible flexible markings is therefore
\begin{equation}\label{word:markedweight}
 u^n(1+t)^{b(M,\mathcal T)}.
\end{equation}

\paragraph{The activity sum.}
We next prove directly that
\begin{equation}\label{word:activity}
 a(M)=\sum_{\mathcal T\text{ spanning tree of }M}
          (1+t)^{b(M,\mathcal T)}.
\end{equation}
Evaluate the same expression on each connected residual map encountered
during edge deletion and contraction. Its weight does not require a
choice of root. Splitting its subset sum according to the
presence of an edge $e$ gives
\begin{equation}\label{word:deletion-contraction}
 a(M)=
 \begin{cases}
 a(M/e)+a(M\setminus e),&e\text{ is neither a bridge nor a loop},\\
 (1+t)a(M/e),&e\text{ is a bridge},\\
 (1+t)a(M\setminus e),&e\text{ is a loop}.
 \end{cases}
\end{equation}
For a nonloop present edge, contraction decreases $|A|$ and $|V|$ by
one and leaves the exponent unchanged. At a bridge, pairing the terms
with and without $e$ gives relative factors $1$ and $t$. At a loop,
including $e$ multiplies the corresponding deleted term by $t$.
The terminal one-vertex map has weight one.

Use the tour on the original darts to select the next edge in this
recursion. When an edge is first reached, a contraction declares it a
tree edge, and a deletion declares it a non-tree edge; thereafter use
the corresponding tour rule whenever that dart is visited. At a bridge
or loop the choice is forced. The next undecided edge depends only on
previous choices. Indeed, contracted edges form a forest, and deletion
of nonbridges preserves a connected residual graph, so every partial
branch has a spanning-tree completion. Its tour agrees with the partial
tour already traversed. Since a completed tree tour visits all darts,
the partial tour cannot close before every edge has been decided.
Consequently the recursion produces each spanning tree exactly once.

Fix a completed branch with spanning tree $\mathcal T$. A tree edge $e$
splits $\mathcal T$ into two components. A non-tree edge joins these two
components precisely when its occurrences alternate with those of $e$:
the contour interval cut out by the pair for $e$ visits exactly one side
of this tree cut. At the decision for $e$, earlier tree contractions
cannot identify the two sides, and the earlier non-tree edges have been
deleted. Hence $e$ is then a bridge if and only if no later non-tree
edge crosses its cut, which is exactly its activity condition.

Similarly, for a non-tree edge $e$, the tree edges crossing it in tour
order are exactly those on the tree path joining its endpoints. At its
decision, it has become a loop if and only if every edge of that path
has already been contracted. This again says exactly that no crossing
edge has a later first occurrence. The forced choices in the branch
are therefore its active edges. Equation~\eqref{word:deletion-contraction}
assigns that branch weight $(1+t)^{b(M,\mathcal T)}$, proving
Equation~\eqref{word:activity}.

There are no hidden symmetry factors in these counts. An automorphism
fixing the root dart and preserving the edge pairing and vertex cyclic
orders fixes every dart, by connectedness. Thus summing
Equation~\eqref{word:markedweight} over rooted map/tree pairs and using
Equation~\eqref{word:activity} proves
Equation~\eqref{word:identity-equation}. For $n=0$, both sides equal one.
\end{proof}

\begin{corollary}[Catalan upper bound]\label{word:catalan}
For every $n\ge0$,
\begin{equation}\label{word:catalan-equation}
 a_nu^n\le \frac{1}{4^n(n+1)}\binom{2n}{n}
          =O\bigl((n+1)^{-3/2}\bigr).
\end{equation}
\end{corollary}

\begin{proof}
Give a burger increment $+1$ and an order increment $-1$. By
Equation~\eqref{word:letterlaw}, these are independent symmetric
increments. An empty-reducing word must have a nonnegative running sum
and terminal sum zero: at every prefix, each order must already have
consumed a distinct preceding burger. Reflection gives exactly
$(n+1)^{-1}\binom{2n}{n}$ such sign sequences of length $2n$.
Apply Proposition~\ref{word:identity}, and then the usual central
binomial estimate.
\end{proof}

\subsection{The first backward survivor}

At a fixed boundary of an independent word, reveal letters backwards,
prepending one letter at each step. Write $S_j$ for the sum of the first
$j$ revealed burger-minus-order increments, with $S_0=0$. Let
\begin{equation}\label{word:J}
 J=\inf\{j\ge1:\text{the revealed suffix has a surviving burger}\}.
\end{equation}
At time $J$, the reduction consists of one burger and some number $K$
of fixed orders of the opposite type. To justify this assertion, note
that before the triggering letter the suffix has no surviving burgers.
That letter must therefore be a burger. By Lemma~\ref{word:reduction},
it survives precisely when none of the suffix's unfilled orders can
consume it; these orders must all be fixed orders of the other type.
In particular,
\begin{equation}\label{word:stopped-sum}
 S_J=1-K.
\end{equation}

This stopping variable is the one studied in Sheffield's inventory
analysis~\cite[Section~3.1]{Sheffield2016}.
His exact identity is $\E(1+K)=1/p$ in the present regime.
We prove the sufficient inequality $\E K\le1/(2p)<1$ directly,
retaining the covariance argument that makes the strictness visible.

\begin{proposition}[Strict drift]\label{word:drift}
The stopping time $J$ is almost surely finite, and
\begin{equation}\label{word:drift-equation}
 \E K\le\frac{1}{2p}<1.
\end{equation}
\end{proposition}

\begin{proof}
We first prove integrability, then exploit the type symmetry in a
stationary word to improve the bound.

\paragraph{A first integrability bound.}
If $S_j=1$, some burger must survive, so $J$ is no larger than the first
time the simple symmetric walk $S$ hits $1$. This hitting time is finite
almost surely. For example, stopping between $-r$ and $1$ gives
probability $r/(r+1)$ of reaching $1$ first, and letting $r$ increase
proves the assertion. Thus $J<\infty$ almost surely.

Stop the walk at the bounded time $J\wedge m$. Since its increments
have mean zero, Equation~\eqref{word:stopped-sum} gives
\begin{equation}\label{word:optional-stopping}
 \E[-S_m;J>m]=\E[1-K;J\le m].
\end{equation}
On $\{J>m\}$ there are no surviving burgers, so $S_m\le0$.
Consequently $\E[K;J\le m]\le\P(J\le m)$. Monotone convergence yields
$\E K\le1$. This uses no finite-mean assumption on $J$.

\paragraph{Typing flexible orders in a stationary word.}
Now take a two-sided independent sequence $(\omega_i)_{i\in\mathbb Z}$
with the letter law in Equation~\eqref{word:letterlaw}. At each boundary
immediately before index $i$, inspect the past until its first surviving
burger appears. The preceding argument and a countable intersection
make this possible almost surely for every $i$. By
Lemma~\ref{word:reduction}, the freshest survivor cannot change when
still older letters are prepended. Let $T_i\in\{+1,-1\}$ be its type
sign, with type 1 positive and type 2 negative.

Define the signed type increment $D_i$ by
\[
 D_i=
 \begin{cases}
 +1,&\omega_i=\mathtt b_1\text{ or }\mathtt o_2,\\
 -1,&\omega_i=\mathtt b_2\text{ or }\mathtt o_1,\\
 -T_i,&\omega_i=\mathtt F.
 \end{cases}
\]
This construction is shift-covariant, so $(D_i)$ is stationary. Each
$D_i$ is measurable with respect to raw letters at indices at most $i$.
Moreover, its flexible typing agrees with every flexible match internal
to any finite word: older supply cannot alter that match, by
Lemma~\ref{word:reduction}. Hence matched pairs cancel when signed
increments are summed.

Work at boundary $0$, and use $J,K$ for its backward variables. Set
\[
 T=T_0,\qquad D^{(m)}=\sum_{i=-m}^{-1}D_i.
\]
On $\{J=j\le m\}$ the final $j$ letters have signed sum $T(1+K)$:
the one burger and the $K$ opposite-type orders all contribute sign
$T$. The earlier signed sum $\sum_{i=-m}^{-j-1}D_i$ depends only on
raw letters earlier than $-j$. It is independent of the final $j$
letters, which determine $\{J=j\}$ and $T$ on that event.
Interchanging the two types preserves this event and reverses $T$,
so $\E[T\mathbf1_{\{J=j\}}]=0$. The earlier signed sum therefore
contributes zero in expectation against $T\mathbf1_{\{J=j\}}$.

On $\{J>m\}$ the reduced suffix consists only of orders. Internal
matched pairs cancel, so $|D^{(m)}|$ is at most the number of unfilled
orders, namely $-S_m$. Combining these observations with
Equation~\eqref{word:optional-stopping} gives
\begin{align}
 \E[T D^{(m)}]
 &\ge \E[1+K;J\le m]-\E[-S_m;J>m]\notag\\
 &=2\E[K;J\le m].\label{word:covariance}
\end{align}
Every variable in this calculation is integrable for fixed $m$; in
particular, $K\le m$ on $\{J\le m\}$.

\paragraph{Nonnegative second moments force strict drift.}
Let $V_m=\E[(D^{(m)})^2]$. Conditional on the raw past before index $0$,
the nonflexible letters have mean signed contribution zero, while the
flexible contribution has conditional mean $-(p/2)T$. Since $D_0^2=1$
and stationarity identifies the second moment of
$\sum_{i=-m}^{0}D_i$ with $V_{m+1}$, we obtain
\begin{equation}\label{word:variance-increment}
 V_{m+1}-V_m
   =1-p\E[T D^{(m)}]
   \le 1-2p\E[K;J\le m].
\end{equation}
If $2p\E K>1$, the last expression is bounded above by a strictly
negative constant for all sufficiently large $m$, by monotone
convergence. Summing Equation~\eqref{word:variance-increment} would then
force some $V_m$ to be negative. This contradiction proves
$\E K\le1/(2p)$, and $p>1/2$ makes the bound strictly smaller than one.
\end{proof}
